\section{Introduction}
\label{sec:introduction}

Court mandates can correct nondelivery while changing the state's production problem. When a judge orders a public agency to provide a medicine, the relevant choice is no longer only whether the patient receives treatment. The agency must also source the good under legal urgency, with accountability concentrated on nondelivery. A procurement office facing short deadlines and one-sided sanctions cannot always wait for aggregation, batching, or regular supplier competition. The public-economics question is therefore which procurement margin the state gives up when delivery is compelled under sanctions.

Higher prices under judicial urgency have two distinct interpretations. They may reflect same-firm markups: suppliers see that a buyer is sanctioned and charge more for the same item. They may instead reflect sourcing reallocation: the state buys smaller lots, loses scale, attracts fewer bidders, or matches with different suppliers. Standard procurement comparisons have trouble separating these mechanisms because prices, quantities, suppliers, and urgency move together. A court-mandated purchase can look expensive even if no incumbent supplier raises its price conditional on selling the same item to the same buyer.

This distinction connects several literatures. Research on public procurement documents large price dispersion, passive waste, and the role of buyer capacity \citep{bandiera2009active,best2023procurement,bosio2022public,decarolis2018buyer}; work on bureaucracy and accountability shows that public agents respond to one-sided pressure and distorted incentives \citep{prendergast2007bureaucrats,bandiera2021allocation}; and work on court efficiency studies how courts shape procurement performance \citep{coviello2018court}. A separate literature on right-to-health litigation studies access, priority setting, and the fiscal pressure of judicialized health claims \citep{ferraz2009right,wang2015right,biehl2009will}. We link these by opening a different court-procurement channel: court-mandated urgent delivery under one-sided compliance pressure. The object is not to restate the aggregate fiscal pressure of litigation, but to identify whether the procurement-cost margin comes from pricing by the same suppliers or from the sourcing process that legal urgency creates.

The setting is S\~ao Paulo pharmaceutical procurement. The state buys medicines through BEC, an electronic procurement platform that records items, prices, quantities, buyers, suppliers, modalities, and tender outcomes. Right-to-health litigation creates urgent court-mandated purchases. Court orders often impose short delivery deadlines and penalties for noncompliance; procurement notices may reveal the judicial order or the administrative process generated by it. The sanction is asymmetric in the procurement market. Officials must deliver or face legal consequences, while suppliers remain voluntary bidders who choose whether and how to participate.

S\~ao Paulo also has an administrative request channel. Some patients obtain medicines outside litigation, generating urgent procurement without judicial sanctions. This comparison is valuable not because it randomizes sanction exposure, but because it is the closest feasible urgent-procurement comparison. Administrative urgent purchases share the compressed execution environment, yet they are selected and larger. The empirical design therefore combines three pieces of evidence: Lee bounds for selection, within firm-buyer-item comparisons for same-firm pricing, and direct sourcing tests for scale and supplier-set reallocation.

We first document the broader urgent-procurement margin. Relative to ordinary purchases of comparable items, urgent purchases have \BPnegPctHeadline{} higher negotiated prices, \BPfirmsPctHeadlineAbs{} fewer bidders, and a \BPsuccessPP{} higher tender-success rate. These estimates establish the procurement environment rather than the paper's sanction design: urgency raises costs and weakens competition even as the state completes tenders more often.

The under-the-gun comparison is within urgent procurement. Court-mandated purchases are costlier than administrative urgent purchases, with a naive litigated-over-administrative gap of \BPutgPointNaive{}. Because administrative requests are screened and selected, we trim the overrepresented administrative group within item-year-PBU strata. The Lee bounds place the selection-corrected litigated-over-administrative gap between \BPutgBoundLow{} and \BPutgBoundHigh{}, and wild cluster inference supports the preferred estimate ($p=\BPutgBoottestPval{}$). The design does not claim random assignment to sanctions. It asks how much of the observed gap remains after disciplining the most direct urgent comparison for selection.

The pricing test then asks whether the same firm charges a sanctioned buyer more for the same item. In firm-buyer-item triples observed under both urgent regimes, the administrative coefficient is \BPutgTripleCoefUTG{} with standard error \BPutgTripleSEUTG{}. Deep markets are repeated urgent settings with enough scale or standardized demand to support within-triple comparisons, mainly above-median quantity and SUS-formulary purchases. In those subsamples the same-firm markup is absent: the corresponding coefficients are \BPwfRobAboveQtyCoef{} and \BPwfRobSUSCoef{}. The result is not a claim that supplier leverage never matters. In below-median quantity purchases and in the earlier period, the coefficients are \BPwfRobBelowQtyCoef{} and \BPwfRobEarlyCoef{}. In deep repeated urgent markets, the sanction-related cost margin does not appear as a broad same-firm markup; in thinner markets, the markup channel can reappear.

The sourcing evidence shows where the margin operates. Administrative urgent orders are \BPutgQtyAdminFoldPref{} times larger than litigated urgent orders, so fragmented court-mandated buying gives up scale. Figure~\ref{fig:sourcing_pricing} decomposes the observed administrative-minus-litigated price gap into a mechanical quantity component, a within-firm component, and a residual composition component. The residual is not, by itself, proof of sourcing. The direct evidence is supplier reallocation: among \BPwinnerSwitchIBpairs{} item-buyer pairs observed under both urgent regimes, mean winner-set Jaccard similarity is \BPwinnerSwitchJaccardMean{}, \BPwinnerSwitchNoOverlapPct{} have no winner overlap, and the modal winner differs in \BPwinnerSwitchDiffModalPct{} of pairs. Conditional on the same firm, the broad markup is weak in deep markets; unconditionally, the state often buys from a different supplier set.

The contribution is to identify the procurement mechanism behind a legal-enforcement margin. The paper separates same-firm pricing from supplier-set sourcing in a setting where standard procurement data make the two observationally equivalent, and it traces how right-to-health litigation changes the buying process that links a court order to a procurement-cost outcome. The pattern is a judicial-enforcement form of passive waste: one-sided sanctions secure delivery but weaken demand aggregation and supplier matching. The implication is not to weaken access to medicines, but to design procurement institutions that preserve aggregation under legal urgency.

The paper proceeds as follows. Section~\ref{sec:institutional} describes the setting, Section~\ref{sec:data} explains the data and samples, Section~\ref{sec:framework} sets out the empirical comparisons, Sections~\ref{sec:results} and~\ref{sec:robustness} report results and robustness, and Section~\ref{sec:policy} discusses interpretation and limits.
